% Figure from MATH 551 notes, section 18. Compile with the notes preamble.
\begin{tikzpicture}
\begin{axis}[nfig, width=0.66\textwidth, height=0.50\textwidth, clip=false,
  axis x line=middle, xmin=0, xmax=1.1, ymin=-1.2, ymax=4.4,
  xtick={0.25,0.5,0.75,1}, xticklabels={$\frac14$,$\frac12$,$\frac34$,$1$},
  ytick={-1,1,2,3,4}, xlabel={}]
  % g = total variation function T(x)
  \addplot[figB, very thick, domain=0:0.25, samples=40] {sin(deg(2*pi*x))};
  \addplot[figB, very thick, domain=0.25:0.75, samples=60] {2-sin(deg(2*pi*x))};
  \addplot[figB, very thick, domain=0.75:1, samples=40] {4+sin(deg(2*pi*x))};
  % h = T - f
  \addplot[figR, very thick, domain=0:0.25, samples=2] {0};
  \addplot[figR, very thick, domain=0.25:0.75, samples=60] {2-2*sin(deg(2*pi*x))};
  \addplot[figR, very thick, domain=0.75:1, samples=2] {4};
  % f
  \addplot[black, very thick, domain=0:1, samples=120] {sin(deg(2*pi*x))};
  \draw[black!35, densely dotted] (axis cs:0.25,0) -- (axis cs:0.25,1);
  \draw[black!35, densely dotted] (axis cs:0.75,-1) -- (axis cs:0.75,3);
  \node[font=\small, anchor=south east] at (axis cs:0.1,0.6) {$f$};
  \node[font=\small, figB, anchor=north west] at (axis cs:0.86,3.3) {$g$};
  \node[font=\small, figR, anchor=south] at (axis cs:0.87,4.02) {$h$};
\end{axis}
\end{tikzpicture}
